\section{神经网络依存分析器}

\subsection{从离散特征到分布式表示}

\begin{figure}[H]
\centering
\includegraphics[width=0.95\textwidth]{slides-images/slide-034.jpg}
\caption{传统特征方法的问题：稀疏、不完备、计算昂贵\protect\footnotemark}
\end{figure}
\footnotetext{Slides 第34页。}

针对传统特征方法的三大问题，自然的解决方案是使用\textbf{分布式表示}（distributed representations）——即我们在前几讲中学习的\textbf{词向量}（word embeddings）。

\subsection{Chen \& Manning (2014) 神经依存分析器}

\begin{figure}[H]
\centering
\includegraphics[width=0.95\textwidth]{slides-images/slide-035.jpg}
\caption{Chen \& Manning (2014) 神经依存分析器架构\protect\footnotemark}
\end{figure}
\footnotetext{Slides 第35页。}

Danqi Chen（当时是 Manning 的博士生，也是 CS224N 的两届 Head TA）构建了第一个成功的\textbf{神经转移依存分析器}，其核心设计如下：

\begin{importantbox}{神经依存分析器的架构}
\textbf{输入表示}：对于当前分析状态，提取以下关键元素的分布式表示并拼接（concatenate）：
\begin{itemize}
    \item 栈顶和缓冲区顶部的\textbf{词向量}（word embeddings）
    \item 对应位置的\textbf{词性向量}（POS tag embeddings）
    \item 已构建的依存弧的\textbf{关系标签向量}（dependency label embeddings）
\end{itemize}

\textbf{网络结构}：
$$h = \text{ReLU}(W_1 \cdot [\text{concatenated inputs}] + b_1)$$
$$\text{output} = \text{softmax}(W_2 \cdot h + b_2)$$

\textbf{输出}：一个概率分布，覆盖所有可能的操作（SHIFT、各种标签的 LEFT-ARC、各种标签的 RIGHT-ARC）。
\end{importantbox}

\begin{figure}[H]
\centering
\includegraphics[width=0.95\textwidth]{slides-images/slide-036.jpg}
\caption{神经网络架构细节：输入层（拼接的分布式表示） $\rightarrow$ 隐藏层（ReLU） $\rightarrow$ softmax 输出层\protect\footnotemark}
\end{figure}
\footnotetext{Slides 第36页。}

\subsection{三种分布式表示的作用}

\begin{knowledgebox}{为什么需要三种嵌入}
\begin{itemize}
    \item \textbf{词向量}（Word Embeddings）：捕获词汇语义相似性。即使训练数据中没有见过某个词在特定位置出现，如果该词与已见过的词语义相近，系统也能做出合理预测
    \item \textbf{词性向量}（POS Tag Embeddings）：将细粒度的词性标签（如单数名词 NN vs. 复数名词 NNS）映射到连续空间，使相似词性获得相近表示
    \item \textbf{依存标签向量}（Dependency Label Embeddings）：捕获已构建的依存关系类型之间的相似性，为后续决策提供结构化上下文
\end{itemize}
\end{knowledgebox}

\subsection{非线性分类器的优势}

传统方法使用\textbf{线性分类器}（逻辑回归、SVM），而神经方法使用\textbf{非线性分类器}（含 ReLU 隐藏层的前馈网络）。这带来了两大优势：

\begin{enumerate}
    \item \textbf{自动学习特征组合}：隐藏层能自动学习输入特征的非线性组合，不再需要手工设计特征交叉
    \item \textbf{更强的表达能力}：非线性分类器能拟合更复杂的决策边界
\end{enumerate}

\subsection{实验结果}

\begin{figure}[H]
\centering
\includegraphics[width=0.95\textwidth]{slides-images/slide-037.jpg}
\caption{实验对比：Chen \& Manning 神经分析器 vs. 传统方法和图方法\protect\footnotemark}
\end{figure}
\footnotetext{Slides 第37页。}

实验结果令人瞩目：

\begin{table}[H]
\centering
\begin{tabular}{lccc}
\toprule
\textbf{分析器} & \textbf{UAS (\%)} & \textbf{LAS (\%)} & \textbf{速度} \\
\midrule
MaltParser（Nivre, 转移+符号特征） & 89.9 & -- & 快 \\
MSTParser（图方法+符号特征） & 91.5 & -- & 慢（$\sim$50x） \\
Chen \& Manning（转移+神经网络） & 92.0 & -- & 快 \\
\bottomrule
\end{tabular}
\caption{不同依存分析器的精度与速度对比}
\end{table}

\begin{importantbox}{神经方法的双重突破}
Chen \& Manning (2014) 的神经依存分析器同时实现了：
\begin{enumerate}
    \item \textbf{精度提升}：超越了传统的转移分析器（MaltParser），达到了图方法（MSTParser）的水平
    \item \textbf{速度保持}：保持了转移分析的线性时间复杂度，甚至比传统转移分析器更快（因为省去了大量特征计算开销）
\end{enumerate}
这打破了精度-速度的传统权衡。
\end{importantbox}

\subsection{本章小结}

Chen \& Manning (2014) 的神经依存分析器用分布式表示替代离散特征、用非线性网络替代线性分类器，在保持线性时间复杂度的同时，将分析精度提升到图方法的水平。三种嵌入（词、词性、依存标签）共同提供丰富的上下文信息。这项工作开创了神经句法分析的新时代。

%% ========== 评估方法 ==========

